\documentclass[11pt,a4paper]{article}

% Packages
\usepackage[utf8]{inputenc}   % Encoding
\usepackage[T1]{fontenc}      % Font encoding
\usepackage{lmodern}          % Modern Latin fonts
\usepackage{hyperref}         % Clickable links
\usepackage{geometry}         % Page layout
\usepackage{listings}
\usepackage{ragged2e}         % Justified text
\usepackage{booktabs}         % Nicer tables
\usepackage{amsmath}          % Math symbols
\lstset{basicstyle=\ttfamily, columns=fullflexible}
\geometry{margin=2cm}

% Document metadata
\title{Skills Tracker --- User Guide}
\date{\today}

\begin{document}
\maketitle
\justifying

\section*{What is this?}

The Skills Tracker lets you self-assess skills in the context of specific fields of use --- for example, how well you apply Python in Data Engineering, or SQL in Reporting.
For each combination of skill and field, you rate four dimensions: theoretical knowledge, practical experience, interest in developing further, and your proficiency within that specific field.
All ratings are tracked over time so both you and management can see how they change.

There is \textbf{nothing to install}. The tool runs on a central server; you open it in any web browser.

\section*{How to access it}

The tool is split into two separate web apps so that personal data is only visible to those who need it.
Your system administrator will give you the URL for your role.

\begin{center}
\begin{tabular}{lll}
\toprule
\textbf{Who you are} & \textbf{Which app} & \textbf{Default URL} \\
\midrule
All staff & Employee app & \texttt{http://<server>:8501} \\
HR / managers & Management app & \texttt{http://<server>:8502} \\
\bottomrule
\end{tabular}
\end{center}

Before entering any data, confirm that \textbf{\texttt{prod\_skills.db}} is shown in the database caption
beneath the page title.

\section*{Employee app}

\subsection*{First time only --- register}

Go to the \textbf{Home} tab, enter your company email address, first name, and last name, then click
\textbf{Continue}. This registers your identity and must be done before using any other tab.

\subsection*{Viewing your skills}

\begin{description}
\item[My Skills] Your current skill profile as a table. Each row is one skill-in-context: the field of use,
  the skill name, and your four ratings for that combination.
  The same skill can appear more than once if you have rated it in different fields.
\item[My Progression] Select a field and a skill to see a line chart of how your four ratings have changed
  over time for that specific combination.
\end{description}

\subsection*{Logging or updating a skill}

Go to the \textbf{Add / Update Skills} tab:
\begin{enumerate}
  \item Select a \textbf{Field of Use} --- the area or context in which you are applying the skill
    (e.g.\ Data Engineering, Reporting, Process Safety). Choose \textbf{+ Add new\ldots} if yours is not listed.
  \item Select the \textbf{Skill}, or choose \textbf{+ Add new\ldots} to add one.
  \item Adjust the four sliders. All four ratings relate to this specific skill \emph{in the chosen field}:
    \begin{itemize}
      \item \emph{Theoretical knowledge} --- how well you understand the underlying concepts (1--5).
      \item \emph{Practical experience} --- how regularly you apply this skill day-to-day (1--5).
      \item \emph{Interest / motivation} --- how keen you are to develop this skill further (1--5).
      \item \emph{Proficiency in [field]} --- your overall proficiency applying this skill within the chosen field (1--5).
    \end{itemize}
  \item Click \textbf{Save}.
\end{enumerate}

Rating Python in Data Engineering and rating Python in Reporting are treated as two separate entries.
Each combination has its own independent history, so changing one does not affect the other.

\section*{Management app}

\begin{description}
\item[Management] Every person's self-assessed values in a filterable table.
  Use the \emph{Filter by skill} and \emph{Filter by person} dropdowns to drill down.
\item[Analytics] Company-wide aggregated charts and skill-gap analysis, filterable by field of use.
  \begin{itemize}
    \item \textbf{Skill Overview table} --- average theoretical, practical, and interest levels per skill,
      plus a knowledge gap column (Theoretical $-$ Practical) and a motivation gap column (Interest $-$ Practical).
    \item \textbf{Grouped bar chart} --- average levels by dimension for each skill.
    \item \textbf{Knowledge Gap bar chart} --- ranked view of where theoretical understanding outpaces practical ability.
    \item \textbf{Gap snapshot tables} --- top knowledge gaps, top motivation gaps, and lowest practical capability.
    \item \textbf{Download reports} --- export the summary table as \texttt{.csv} or \texttt{.xlsx}.
  \end{itemize}
\item[Admin] Enter an email address and, optionally, a skill name, then click \textbf{Delete}.
  Leaving the skill field blank removes \emph{all} skills for that person.
  \textbf{Deletions are permanent and cannot be undone.}
\end{description}

\section*{Important notes}

\begin{itemize}
  \item All colleagues share a single database on the server. Changes made by one person are immediately
    visible to others.
  \item Do not attempt to access, copy, or edit the database file directly --- this risks data corruption.
  \item Deletions via the Admin tab are \textbf{permanent and cannot be undone}.
  \item If the app is unavailable or behaving unexpectedly, contact your system administrator.
    Do not attempt to restart or reinstall anything yourself.
\end{itemize}

\end{document}
